\documentclass[11pt]{article}

\usepackage[a4paper,margin=1in]{geometry}
\usepackage{amsmath,amssymb,amsthm,mathtools}
\usepackage{enumitem}
\usepackage{hyperref}
\usepackage{microtype}
\usepackage{xcolor}

\hypersetup{
  colorlinks=true,
  linkcolor=blue!55!black,
  citecolor=blue!55!black,
  urlcolor=blue!55!black
}

\newtheorem{theorem}{Theorem}[section]
\newtheorem{proposition}[theorem]{Proposition}
\newtheorem{lemma}[theorem]{Lemma}
\newtheorem{corollary}[theorem]{Corollary}
\newtheorem{assumption}[theorem]{Assumption}
\newtheorem{definition}[theorem]{Definition}
\newtheorem{example}[theorem]{Example}
\newtheorem{remark}[theorem]{Remark}

\newcommand{\R}{\mathbb{R}}
\newcommand{\Rp}{\mathbb{R}_{+}}
\newcommand{\Pcal}{\mathcal{P}}
\newcommand{\Bcal}{\mathcal{B}}
\newcommand{\X}{\mathsf{X}}
\newcommand{\Y}{\mathsf{Y}}
\newcommand{\Z}{\mathsf{Z}}
\newcommand{\eps}{\varepsilon}
\newcommand{\dd}{\,\mathrm{d}}
\newcommand{\osc}{\operatorname{osc}}
\newcommand{\TV}{\operatorname{TV}}
\newcommand{\KL}{\operatorname{KL}}
\newcommand{\supp}{\operatorname{supp}}
\newcommand{\argmax}{\operatorname*{arg\,max}}
\newcommand{\argmin}{\operatorname*{arg\,min}}
\newcommand{\one}{\mathbf{1}}
\newcommand{\id}{\mathrm{id}}
\newcommand{\Dom}{\operatorname{dom}}
\newcommand{\Lip}{\operatorname{Lip}}
\newcommand{\esssup}{\operatorname*{ess\,sup}}
\newcommand{\essinf}{\operatorname*{ess\,inf}}
\newcommand{\otimesi}{\mathop{\otimes}}
\newcommand{\defeq}{\coloneqq}

\title{Contraction Mechanisms for Generalized Sinkhorn Maps\\
       under \texorpdfstring{$\phi$}{phi}-Divergence Regularization}
\author{A technical note}
\date{\today}

\begin{document}
\maketitle

\begin{abstract}
Entropic optimal transport has two useful order properties: the soft
$c$-transform is order reversing and commutes with additive constants, hence
the double transform is non-expansive in the variation seminorm.  In the
Kullback--Leibler case, strict contraction follows in compact bounded-cost
settings from Hilbert's projective metric and Birkhoff's contraction theorem;
Carlier's proof gives another route through boundedness of Sinkhorn iterates,
strong convexity of the exponential on bounded intervals, and a
Gauss--Seidel/coordinate-descent comparison.  This note explains which parts
survive when the KL penalty is replaced by a general $\phi$-divergence.

The main message is that monotonicity and non-expansiveness are structural,
but strict contraction is not.  For a smooth interior generalized transform,
its derivative is a Markov averaging operator whose transition probabilities
are proportional to the curvature of the dual integrand
$\Phi=\eps\phi_+^*(\cdot/\eps)$.  Therefore the correct replacement for
positivity of the Gibbs kernel is a uniform overlap, or Doeblin, condition on
these curvature kernels.  A simple sufficient condition is a uniform strip on
which $0<m\leq \Phi''\leq M<\infty$.  The Carlier argument extends by replacing
the exponential inequalities with the corresponding strong-convexity and
Lipschitz estimates for $\Phi$ on that strip.  Without such an interior
curvature condition, sparse divergences such as quadratic/chi-square
regularization can have local contraction factor one.
\end{abstract}

\tableofcontents

\section{Introduction}

Let $(\X,\alpha)$ and $(\Y,\beta)$ be probability spaces and let
$c:\X\times\Y\to\R$ be a measurable cost.  For a convex function
$\phi:[0,\infty)\to(-\infty,+\infty]$ with $\phi(1)=0$, the
$\phi$-regularized transport problem is
\begin{equation}\label{eq:primal}
  \inf_{\pi\in\Pi(\alpha,\beta)}
  \left\{
  \int_{\X\times\Y} c\,\dd\pi
  +\eps D_\phi(\pi\mid \alpha\otimes\beta)
  \right\},
  \qquad \eps>0.
\end{equation}
When $\phi(r)=r\log r-r+1$, this is the usual entropy-regularized optimal
transport problem.  Replacing KL by a general $\phi$ leads to a generalized
Sinkhorn algorithm: one alternately maximizes the dual in one potential while
keeping the other fixed.  This generalized transform was introduced and
studied in the $f$-divergence OT literature; in particular, Terjek and
Gonzalez-Sanchez formulate the $(c,\eps,\phi)$-transform, its implicit equation,
and convergence conditions for the generalized Sinkhorn algorithm
\cite{terjek2022}.  Earlier convergence results for general convex
regularization also appear in Di Marino--Gerolin \cite{dimarino2020}.

For KL, several contraction analyses are available.  The classical two-marginal
compact proof uses Hilbert's projective metric and Birkhoff--Hopf contraction
for positive kernel operators; Carlier instead proves linear convergence of the
multi-marginal Sinkhorn algorithm by combining boundedness of the iterates,
strong convexity and Lipschitz continuity of the exponential on bounded
intervals, and coordinate-descent estimates \cite{carlier2022}.  Recent work
extends Hilbert-metric methods for KL to certain unbounded settings by changing
the underlying cone and allowing kernels that are not bounded away from zero
\cite{eckstein2025}.  These facts suggest the question studied here:

\begin{quote}
\emph{Which contraction mechanisms extend from KL to a general
$\phi$-divergence?}
\end{quote}

The answer is nuanced.  The order-theoretic non-expansiveness extends almost
verbatim.  Strict contraction requires an additional quantitative interiority
condition.  The most transparent way to see this is to differentiate the
generalized transform.  Its derivative is a Markov operator with kernel
proportional to $\Phi''$, where
\begin{equation}\label{eq:Phi}
  \Phi(s) \defeq \eps\,\phi_+^*(s/\eps),
  \qquad
  \phi_+(r)\defeq \phi(r)+\iota_{[0,\infty)}(r).
\end{equation}
Thus strict variation contraction is equivalent, at the infinitesimal level, to
a Dobrushin overlap condition for those curvature kernels.  This is automatic
for KL on compact bounded-cost problems, but it can fail for other divergences.
A two-point chi-square example in Section~\ref{sec:counterexample} shows that
one cannot hope for a theorem saying that all generalized Sinkhorn maps are
strict contractions.

The note is organized as follows.  Section~\ref{sec:dual} recalls the dual and
the generalized soft $c$-transform.  Section~\ref{sec:order} proves order
reversal, translation covariance and non-expansiveness in the variation
seminorm.  Section~\ref{sec:derivative} derives the derivative kernel and the
Dobrushin contraction criterion.  Section~\ref{sec:kl} explains why the KL case
has additional Hilbert-metric structure.  Section~\ref{sec:carlier} formulates
the Carlier-style extension through local strong convexity and smoothness of
$\Phi$.  Section~\ref{sec:counterexample} gives the obstruction.  The final
section sketches weighted/noncompact extensions.

\section{Dual problem and generalized soft \texorpdfstring{$c$}{c}-transform}\label{sec:dual}

\subsection{Normalization of the divergence}

The affine part of $\phi$ does not affect the value of
$D_\phi(\pi\mid \alpha\otimes\beta)$ on probability couplings: adding
$a(r-1)$ to $\phi(r)$ integrates to zero when $\pi$ and
$\alpha\otimes\beta$ have the same total mass.  We therefore use the following
normalization throughout the note.

\begin{assumption}[standing convex-analytic assumptions]\label{ass:phi}
The function $\phi:[0,\infty)\to(-\infty,+\infty]$ is proper, lower
semicontinuous and convex, with
\begin{equation*}
  \phi(1)=0,\qquad 0\in\partial \phi(1).
\end{equation*}
Let $\phi_+=\phi+\iota_{[0,\infty)}$ and let
$\Phi(s)=\eps\phi_+^*(s/\eps)$.  On intervals where differentiability is
needed, assume $\Phi\in C^2(I)$ and write
\begin{equation*}
  \vartheta(s)\defeq \Phi'(s)=(\phi_+^*)'(s/\eps).
\end{equation*}
The normalization gives $\vartheta(0)=1$ whenever $\vartheta$ is defined at
zero.
\end{assumption}

The usual generalized $f$-divergence includes a singular part:
\begin{equation*}
  D_\phi(\pi\mid m)
  =\int \phi\!\left(\frac{\dd\pi}{\dd m}\right)\dd m
  +\phi'_\infty\,\pi^\perp(\X\times\Y),
  \qquad m=\alpha\otimes\beta,
\end{equation*}
where $\pi=\frac{\dd\pi}{\dd m}m+\pi^\perp$ is the Lebesgue decomposition and
$\phi'_\infty=\lim_{r\to\infty}\phi(r)/r$.  If $\phi'_\infty=+
\infty$, singular mass is forbidden.  If $\phi'_\infty<\infty$, the dual has a
corresponding upper constraint.  This is one reason why non-superlinear
regularizations may partially collapse to the ordinary hard $c$-transform
\cite{terjek2022}.

\subsection{Fenchel dual}

Formally, and rigorously under the usual compactness/continuity or integrable
lower-bounded hypotheses, the dual of \eqref{eq:primal} is
\begin{equation}\label{eq:dual}
  \sup_{f,g}
  \left\{
  \int_\X f\,\dd\alpha+
  \int_\Y g\,\dd\beta
  -\int_{\X\times\Y}\Phi\big(f(x)+g(y)-c(x,y)\big)
  \dd\alpha(x)\dd\beta(y)
  \right\}.
\end{equation}
Equivalently, if one does not absorb the effective domain of $\phi_+^*$ into
$\Phi$, then $f(x)+g(y)-c(x,y)\leq \eps\phi'_\infty$ is an explicit constraint.
For KL, $\phi(r)=r\log r-r+1$ and
\begin{equation*}
  \phi_+^*(t)=e^t-1,
  \qquad
  \Phi(s)=\eps(e^{s/\eps}-1),
  \qquad
  \vartheta(s)=e^{s/\eps}.
\end{equation*}
The additive constant $-\eps$ in $\Phi$ is irrelevant for optimization but
convenient for normalization.

\subsection{The generalized transform}

Fix a potential $f$ on $\X$.  The update of the $\Y$-potential is the pointwise
maximizer of the dual in $g(y)$:
\begin{equation}\label{eq:Tdef}
  T_\alpha f(y)
  \defeq
  \argmax_{\gamma}
  \left\{
  \gamma-
  \int_\X \Phi\big(f(x)+\gamma-c(x,y)\big)\dd\alpha(x)
  \right\}.
\end{equation}
When the maximizer is interior and unique, it is characterized by
\begin{equation}\label{eq:implicit}
  \int_\X
  \vartheta\big(f(x)+T_\alpha f(y)-c(x,y)\big)
  \dd\alpha(x)=1.
\end{equation}
This is the generalized soft $c$-transform.  The analogous map
$T_\beta$ sends functions on $\Y$ to functions on $\X$:
\begin{equation}\label{eq:Tbeta}
  \int_\Y
  \vartheta\big(T_\beta g(x)+g(y)-c(x,y)\big)
  \dd\beta(y)=1.
\end{equation}
The generalized Sinkhorn map on $\X$-potentials is
\begin{equation}\label{eq:double}
  S\defeq T_\beta T_\alpha.
\end{equation}

\begin{remark}[boundary cases]
If $\Phi$ is finite only on a half-line or if $\vartheta$ has a finite plateau,
Equation~\eqref{eq:implicit} may fail to have an interior solution.  Then the
maximizer in \eqref{eq:Tdef} occurs at the boundary and the update becomes a
hard $c$-transform plus the endpoint constant.  Terjek--Gonzalez-Sanchez call
out precisely this phenomenon in the non-superlinear case \cite{terjek2022}.
All differential contraction statements below are made on the interior region
where \eqref{eq:implicit} is valid.
\end{remark}

\section{Order structure and variation non-expansiveness}\label{sec:order}

For a bounded measurable function $h$ on a probability space, define the
variation seminorm
\begin{equation*}
  \|h\|_{\mathrm{var}}
  \defeq \osc(h)
  \defeq \esssup h-\essinf h.
\end{equation*}
This is the quotient norm on $L^\infty/\R$.

\begin{proposition}[order reversal and translations]\label{prop:order}
Assume the transform $T_\alpha$ is single-valued.  Then
\begin{align}
  f\leq \tilde f
  &\quad\Longrightarrow\quad
  T_\alpha \tilde f\leq T_\alpha f,
  \label{eq:orderreversal}\\
  T_\alpha(f+a)&=T_\alpha f-a,
  \qquad a\in\R.\label{eq:translation}
\end{align}
Consequently $S=T_\beta T_\alpha$ is order preserving and satisfies
$S(f+a)=Sf+a$.
\end{proposition}

\begin{proof}
The translation identity follows directly from \eqref{eq:Tdef} or
\eqref{eq:implicit}.  For order reversal, suppose $f\leq \tilde f$.  For a fixed
$y$, the function
\begin{equation*}
  \gamma\mapsto
  \int_\X \vartheta(f(x)+\gamma-c(x,y))\dd\alpha(x)
\end{equation*}
is nondecreasing in both $f$ and $\gamma$.  If
$\gamma_f=T_\alpha f(y)$ solves \eqref{eq:implicit}, then at the same value
$\gamma_f$ the corresponding integral for $\tilde f$ is at least one.  To make
it equal to one, the value of $\gamma$ must be no larger.  The same argument can
be written at the level of the concave maximization problem \eqref{eq:Tdef}; in
nonsmooth cases one uses monotonicity of the subdifferential of $\Phi$.
\end{proof}

\begin{proposition}[non-expansiveness in variation]\label{prop:nonexp}
Let $T$ be an order-reversing map with $T(f+a)=Tf-a$.  Then
\begin{equation}\label{eq:Tnonexp}
  \osc(Tf-T\tilde f)\leq \osc(f-\tilde f).
\end{equation}
If $S$ is order preserving and $S(f+a)=Sf+a$, then
\begin{equation}\label{eq:Snonexp}
  \osc(Sf-S\tilde f)\leq \osc(f-\tilde f).
\end{equation}
In particular, the generalized double transform is non-expansive in
$L^\infty/\R$ whenever the transforms are single-valued.
\end{proposition}

\begin{proof}
Let
\begin{equation*}
  m\defeq \essinf(f-\tilde f),
  \qquad
  M\defeq \esssup(f-\tilde f).
\end{equation*}
Then $\tilde f+m\leq f\leq \tilde f+M$.  Applying an order-reversing $T$ and
using translation covariance gives
\begin{equation*}
  T\tilde f-M\leq Tf\leq T\tilde f-m.
\end{equation*}
Therefore $Tf-T\tilde f\in[-M,-m]$ a.e., proving \eqref{eq:Tnonexp}.  The
order-preserving case is identical, with
$S\tilde f+m\leq Sf\leq S\tilde f+M$.
\end{proof}

\begin{remark}
This argument gives non-expansiveness only.  It contains no mechanism forcing a
strict improvement of the oscillation.  Strictness depends on how much averaging
is performed by the transform.  The next section identifies the averaging
kernel.
\end{remark}

\section{Derivative kernels and Dobrushin contraction}\label{sec:derivative}

\subsection{Derivative of a generalized transform}

Assume for this section that the update is interior, that $\Phi\in C^2$ on the
range of all quantities involved, and that
\begin{equation*}
  \int_\X \Phi''\big(f(x)+T_\alpha f(y)-c(x,y)\big)\dd\alpha(x)>0
\end{equation*}
for every $y$.  Put
\begin{equation}\label{eq:zy}
  z_y^f(x)
  \defeq f(x)+T_\alpha f(y)-c(x,y).
\end{equation}

\begin{proposition}[linearization]\label{prop:derivative}
The Frechet derivative of $T_\alpha$ at $f$ is
\begin{equation}\label{eq:derivative}
  D T_\alpha(f)[h](y)
  =-
  \int_\X h(x)\,P_y^f(\dd x),
\end{equation}
where $P_y^f$ is the probability measure
\begin{equation}\label{eq:curvaturekernel}
  P_y^f(\dd x)
  =p_y^f(x)\alpha(\dd x),
  \qquad
  p_y^f(x)
  =
  \frac{\Phi''(z_y^f(x))}
       {\int_\X \Phi''(z_y^f(x'))\dd\alpha(x')}.
\end{equation}
\end{proposition}

\begin{proof}
Differentiate the identity
\begin{equation*}
  \int_\X \Phi'\big(f(x)+T_\alpha f(y)-c(x,y)\big)\dd\alpha(x)=1
\end{equation*}
in the direction $h$.  This gives
\begin{equation*}
  \int_\X \Phi''(z_y^f(x))
  \big(h(x)+D T_\alpha(f)[h](y)\big)\dd\alpha(x)=0,
\end{equation*}
which is exactly \eqref{eq:derivative}.
\end{proof}

Thus the derivative of $T_\alpha$ is the negative of a Markov averaging
operator.  The minus sign is harmless for oscillations.

\subsection{Dobrushin coefficient}

For a Markov kernel $P=(P_y)_{y\in\Y}$ from $\Y$ to $\X$, define its Dobrushin
coefficient
\begin{equation}\label{eq:dobrushin}
  \delta(P)
  \defeq
  \sup_{y,y'} \|P_y-P_{y'}\|_{\TV}
  =1-\inf_{y,y'}\int_\X \min(p_y,p_{y'})\dd\alpha,
\end{equation}
whenever $P_y=p_y\alpha$.  Here
$\|\mu-\nu\|_{\TV}=\sup_A|\mu(A)-\nu(A)|=\frac12\int |p-q|\dd\alpha$.
The standard Dobrushin estimate says
\begin{equation}\label{eq:dobrushin-est}
  \osc\left(y\mapsto \int h\dd P_y\right)
  \leq \delta(P)\,\osc(h).
\end{equation}

Combining this estimate with Proposition~\ref{prop:derivative} gives the
infinitesimal contraction factor of the generalized transform.

\begin{corollary}[local variation contraction]\label{cor:local}
At an interior smooth point $f$,
\begin{equation}\label{eq:localfactor}
  \osc(DT_\alpha(f)[h])
  \leq \delta(P^f)\,\osc(h),
\end{equation}
where $P^f$ is the curvature kernel \eqref{eq:curvaturekernel}.  In particular,
$T_\alpha$ is strictly contracting at $f$ only if the family
$(P_y^f)_{y\in\Y}$ has nontrivial overlap:
\begin{equation}\label{eq:overlap}
  \lambda(f)
  \defeq
  \inf_{y,y'}\int_\X \min(p_y^f,p_{y'}^f)\dd\alpha>0.
\end{equation}
The corresponding infinitesimal factor is at most $1-\lambda(f)$.
\end{corollary}

\subsection{A nonlinear contraction criterion}

The previous statement is infinitesimal.  A uniform estimate along line
segments gives a genuine Lipschitz contraction.

\begin{theorem}[Dobrushin criterion for the generalized transform]\label{thm:dobrushin}
Let $\mathcal{U}$ be a convex class of potentials on $\X$ such that
$T_\alpha$ is interior and $C^1$ on $\mathcal{U}$.  Assume that
\begin{equation}\label{eq:uniform-dob}
  \delta(P^f)\leq \kappa_X<1,
  \qquad f\in\mathcal{U}.
\end{equation}
Then
\begin{equation}\label{eq:Tcontract}
  \osc(T_\alpha f-T_\alpha \tilde f)
  \leq \kappa_X\,\osc(f-\tilde f),
  \qquad f,\tilde f\in\mathcal{U}.
\end{equation}
If, similarly, along the corresponding $\Y$-potentials one has
\begin{equation*}
  \delta(Q^g)\leq \kappa_Y<1
\end{equation*}
for the derivative kernels of $T_\beta$, then the double transform satisfies
\begin{equation}\label{eq:Scontract}
  \osc(Sf-S\tilde f)
  \leq \kappa_X\kappa_Y\,\osc(f-\tilde f).
\end{equation}
\end{theorem}

\begin{proof}
For $f_t=(1-t)f+t\tilde f$,
\begin{equation*}
  T_\alpha\tilde f-T_\alpha f
  =\int_0^1 DT_\alpha(f_t)[\tilde f-f]\dd t.
\end{equation*}
Using \eqref{eq:localfactor} and \eqref{eq:uniform-dob} yields
\eqref{eq:Tcontract}.  The estimate for $S=T_\beta T_\alpha$ follows by applying
the one-block estimate twice.
\end{proof}

\subsection{Uniform curvature strip}

The Dobrushin overlap condition is conceptually exact but may be hard to check.
A crude, useful sufficient condition is uniform curvature on a strip.

\begin{corollary}[curvature strip implies contraction]\label{cor:strip}
Assume that for every potential under consideration and every $x,y$,
\begin{equation}\label{eq:strip}
  z_y^f(x)=f(x)+T_\alpha f(y)-c(x,y)\in I,
\end{equation}
where $I\subset\R$ is an interval on which
\begin{equation}\label{eq:muL}
  0<\mu_I\defeq\inf_{s\in I}\Phi''(s)
  \leq
  L_I\defeq\sup_{s\in I}\Phi''(s)<\infty.
\end{equation}
Then
\begin{equation}\label{eq:stripfactor}
  \osc(T_\alpha f-T_\alpha\tilde f)
  \leq
  \left(1-\frac{\mu_I}{L_I}\right)
  \osc(f-\tilde f)
\end{equation}
provided the same strip condition holds along the segment between $f$ and
$\tilde f$.
\end{corollary}

\begin{proof}
For each $y$,
\begin{equation*}
  \int_\X \Phi''(z_y^f(x))\dd\alpha(x)\leq L_I,
\end{equation*}
while $\Phi''(z_y^f(x))\geq \mu_I$.  Hence the density
$p_y^f$ in \eqref{eq:curvaturekernel} satisfies
$p_y^f(x)\geq \mu_I/L_I$ for $\alpha$-a.e. $x$.  Therefore
\begin{equation*}
  \int \min(p_y^f,p_{y'}^f)\dd\alpha\geq \mu_I/L_I,
\end{equation*}
and Theorem~\ref{thm:dobrushin} applies.
\end{proof}

\begin{remark}[the strip estimate is not sharp]
For KL, $\Phi''(s)=\eps^{-1}e^{s/\eps}$.  On a strip of width $R$ one gets the
rough factor $1-e^{-R/\eps}$.  Hilbert-metric estimates exploit the special
multiplicative structure of the exponential and typically produce sharper or
more intrinsic constants.  For a general $\phi$, however, the curvature-kernel
criterion is the natural replacement.
\end{remark}

\section{What is special about KL?}\label{sec:kl}

For KL,
\begin{equation*}
  \vartheta(s)=e^{s/\eps}.
\end{equation*}
The optimal density associated with dual potentials is
\begin{equation}\label{eq:KLdensity}
  \rho(x,y)
  =e^{(f(x)+g(y)-c(x,y))/\eps}
  =a(x)K(x,y)b(y),
\end{equation}
where
\begin{equation*}
  a(x)=e^{f(x)/\eps},\qquad
  b(y)=e^{g(y)/\eps},\qquad
  K(x,y)=e^{-c(x,y)/\eps}.
\end{equation*}
Sinkhorn updates are therefore positive linear kernel operations followed by
diagonal rescalings.  Hilbert's projective metric is built precisely to study
positive homogeneous maps on cones.  If $K$ is bounded above and below on a
compact domain, Birkhoff--Hopf gives strict projective contraction.  Recent work
extends this strategy for KL by changing the cone and allowing controlled decay
of $K$ in unbounded settings \cite{eckstein2025}.

For a general divergence,
\begin{equation}\label{eq:general-density}
  \rho(x,y)=\vartheta(f(x)+g(y)-c(x,y)).
\end{equation}
There is no factorization of the form $a(x)K(x,y)b(y)$ unless
$\vartheta$ is exponential.  Indeed, a positive function satisfying
$\vartheta(r+s)=A(r)B(s)$, under minimal regularity, is exponential up to
constants.  Thus the direct Hilbert-metric proof is a KL-specific argument.

The derivative formula shows what remains.  Even though the nonlinear map is no
longer a projective positive linear map, its linearization is still a Markov
averaging operator, now with weights proportional to $\Phi''$.  Consequently:
\begin{equation*}
  \text{KL Hilbert positivity}
  \quad\leadsto\quad
  \text{uniform overlap of curvature kernels for general }\phi.
\end{equation*}
This is the main structural replacement.

\section{Carlier-style extension through strong convexity}\label{sec:carlier}

Carlier's proof of linear convergence for entropic multi-marginal Sinkhorn is
based on three ingredients: bounded Sinkhorn iterates, strong convexity of the
exponential on the bounded range reached by the iterates, and a standard
cyclic-coordinate-descent comparison \cite{carlier2022}.  This mechanism is not
specific to the exponential.  It extends to any $\Phi$ satisfying uniform
strong convexity and smoothness on the relevant range.

\subsection{Multi-marginal dual functional}

Let $(\X_i,m_i)$, $i=1,\ldots,N$, be probability spaces and let
$m=m_1\otimes\cdots\otimes m_N$.  Write
$x=(x_1,\ldots,x_N)$ and let $c:\prod_i\X_i\to\R$.  For potentials
$u=(u_1,\ldots,u_N)$, define
\begin{equation}\label{eq:multiF}
  F(u)
  \defeq
  \int_{\prod_i\X_i}
  \Phi\left(\sum_{i=1}^N u_i(x_i)-c(x)\right)
  \dd m(x)
  -
  \sum_{i=1}^N\int_{\X_i} u_i\dd m_i.
\end{equation}
Maximizing the dual is equivalent to minimizing $F$.  The functional is
invariant under adding constants $a_i$ to the blocks whenever
$\sum_i a_i=0$, so one works on a gauge-fixed affine space, for example
\begin{equation*}
  E=\left\{u:\int u_i\dd m_i=0\text{ for }i=1,
  \ldots,N-1\right\}.
\end{equation*}
Cyclic generalized Sinkhorn is exact block minimization of $F$ over
$u_1,u_2,\ldots,u_N$, repeated.

The block optimality condition for the $i$-th update is
\begin{equation}\label{eq:blockcondition}
  \int_{\prod_{j\ne i}\X_j}
  \vartheta\left(
  u_i(x_i)+\sum_{j\ne i}u_j(x_j)-c(x)
  \right)
  \dd m_{-i}(x_{-i})=1
\end{equation}
for $m_i$-a.e. $x_i$, with the other blocks fixed.  For $N=2$ this is exactly
\eqref{eq:implicit} or \eqref{eq:Tbeta}.

\subsection{The two scalar inequalities that replace KL}

Let $I\subset\R$ be an interval and suppose
\begin{equation}\label{eq:scalar-mu-L}
  0<\mu_I\leq \Phi''(s)\leq L_I<\infty,
  \qquad s\in I.
\end{equation}
Then for all $a,b\in I$,
\begin{align}
  \Phi(b)-\Phi(a)-\Phi'(a)(b-a)
  &\geq \frac{\mu_I}{2}(b-a)^2,\label{eq:strongPhi}\\
  |\Phi'(b)-\Phi'(a)|
  &\leq L_I|b-a|.\label{eq:smoothPhi}
\end{align}
In Carlier's KL proof these are used with $\Phi(s)=e^s$ on a bounded interval:
\begin{equation*}
  e^b-e^a-e^a(b-a)\geq \frac{e^{-M}}{2}(b-a)^2,
  \qquad
  |e^b-e^a|\leq e^M|b-a|,
  \qquad a,b\in[-M,M].
\end{equation*}
Thus \eqref{eq:strongPhi}--\eqref{eq:smoothPhi} are the only replacements
needed at the scalar level.

\subsection{A reusable sufficient theorem}

The following theorem is deliberately stated as a sufficient criterion, not as
an optimal convergence theorem.  Its purpose is to isolate what must be checked
for a given divergence and a given class of costs.

\begin{theorem}[Carlier-type linear convergence for general $\phi$]\label{thm:carlier-general}
Consider the functional $F$ in \eqref{eq:multiF} on a gauge-fixed space $E$.
Assume:
\begin{enumerate}[label=(A\arabic*)]
  \item \label{A1} The cyclic block minimizers are well-defined and remain in an
  interior region.
  \item \label{A2} There exists an interval $I$ such that, for every iterate
  $u^t$, every intermediate Gauss--Seidel state in the cycle, and an optimal
  potential $u^\star$,
  \begin{equation}\label{eq:iterates-strip}
    \sum_{i=1}^N u_i(x_i)-c(x)\in I
  \end{equation}
  for $m$-a.e. $x$.
  \item \label{A3} On $I$, $\Phi$ satisfies
  $0<\mu_I\leq \Phi''\leq L_I<\infty$.
  \item \label{A4} On the chosen gauge $E$, the minimizer $u^\star$ is unique.
\end{enumerate}
Then the generalized Sinkhorn/Gauss--Seidel iterates converge to $u^\star$ in
$L^2$ block norm, and the objective values converge linearly.  More precisely,
one obtains a non-optimized estimate of the form
\begin{equation}\label{eq:generic-rate}
  F(u^t)-F(u^\star)
  \leq
  \left(1-\frac{1}{N}\left(\frac{\mu_I}{L_I}\right)^2\right)^t
  \big(F(u^0)-F(u^\star)\big),
\end{equation}
up to harmless changes of the numerical constant depending on the precise gauge
and block norm convention.
\end{theorem}

\begin{proof}[Proof sketch]
Let $u^{t,i}$ be the intermediate state after blocks $1,\ldots,i$ have been
updated during cycle $t$, with $u^{t,0}=u^t$ and $u^{t,N}=u^{t+1}$.  Since block
$i$ is exactly minimized, the first variation in that block vanishes at
$u^{t,i}$.  Applying \eqref{eq:strongPhi} to the change from $u^{t,i-1}$ to
$u^{t,i}$ and summing over the blocks yields the sufficient-decrease estimate
\begin{equation}\label{eq:suff-dec}
  F(u^t)-F(u^{t+1})
  \geq
  \frac{\mu_I}{2}
  \sum_{i=1}^N
  \|u_i^{t+1}-u_i^t\|_{L^2(m_i)}^2 .
\end{equation}
This is the direct analogue of Carlier's Lemma 3.2.

To compare the current gap with the step size, use convexity at $u^t$ and insert
intermediate states $u^{t,i}$ for which the block gradients vanish.  The
difference between the gradient at $u^t$ and the gradient at $u^{t,i}$ is
controlled by \eqref{eq:smoothPhi}; Young's inequality and
\eqref{eq:strongPhi} give an estimate of the type
\begin{equation}\label{eq:gap-step}
  F(u^t)-F(u^\star)
  \leq
  \frac{N L_I^2}{2\mu_I}
  \sum_{i=1}^N
  \|u_i^{t+1}-u_i^t\|_{L^2(m_i)}^2 .
\end{equation}
Combining \eqref{eq:suff-dec} and \eqref{eq:gap-step} gives
\begin{equation*}
  F(u^t)-F(u^{t+1})
  \geq
  \frac{1}{N}\left(\frac{\mu_I}{L_I}\right)^2
  \big(F(u^t)-F(u^\star)\big),
\end{equation*}
which is \eqref{eq:generic-rate}.  The $L^2$ convergence follows from summing
\eqref{eq:suff-dec}, compactness or weak compactness of bounded iterates, and
uniqueness on the gauge-fixed quotient, as in the KL proof.
\end{proof}

\subsection{How to obtain the strip}

For bounded costs and normalized $\vartheta(0)=1$, the transform itself gives
useful oscillation estimates.  Fix $y$ and write
$a_y(x)=f(x)-c(x,y)$.  Equation~\eqref{eq:implicit} is
\begin{equation*}
  \int \vartheta(a_y(x)+\gamma)\dd\alpha(x)=1.
\end{equation*}
Since $\vartheta$ is increasing and $\vartheta(0)=1$, the solution satisfies
\begin{equation}\label{eq:gamma-range}
  -\esssup_x a_y(x)
  \leq \gamma
  \leq
  -\essinf_x a_y(x).
\end{equation}
Equivalently,
\begin{equation}\label{eq:zrange}
  z_y^f(x)=a_y(x)+\gamma
  \in
  [-\osc(a_y),\osc(a_y)]
  \quad \text{a.e.}
\end{equation}
Moreover, comparing two points $y,y'$ gives
\begin{equation}\label{eq:Lip-cost}
  |T_\alpha f(y)-T_\alpha f(y')|
  \leq
  \esssup_x |c(x,y)-c(x,y')|.
\end{equation}
Thus after one block update the oscillation of the updated potential is
controlled by the oscillation of the cost, independently of the incoming
potential.  In bounded-cost settings this usually gives bounded potentials
after fixing a gauge, hence a bounded interval containing the values of
$\sum_i u_i-c$.

This bounded interval is not yet enough: it must lie in a region where
$\Phi''$ is bounded above and below.  KL satisfies this on every bounded
interval.  Many other divergences satisfy it only in a proper interior region.
For sparse regularizers, the lower bound may fail because $\Phi''$ vanishes on
inactive parts of the domain.

\section{A counterexample: chi-square can have contraction factor one}\label{sec:counterexample}

The obstruction is not merely technical.  Some $\phi$-divergences generate
sparse couplings, and sparse active sets can destroy strict contraction.

Consider
\begin{equation*}
  \phi(r)=\frac12(r-1)^2,
  \qquad r\geq 0.
\end{equation*}
Then
\begin{equation*}
  \phi_+^*(t)=
  \begin{cases}
  t+\frac12t^2, & t\geq -1,\\
  -\frac12, & t<-1,
  \end{cases}
  \qquad
  \vartheta(s)=\Phi'(s)=\left(1+\frac{s}{\eps}\right)_+,
\end{equation*}
so
\begin{equation*}
  \Phi''(s)=\frac1\eps\one_{\{s>-\eps\}}
\end{equation*}
away from the kink.  The curvature kernel is uniform on the active set
$\{s>-\eps\}$ and zero outside it.

Let
\begin{equation*}
  \X=\Y=\{1,2\},
  \qquad
  \alpha=\beta=\left(\frac12,\frac12\right),
\end{equation*}
and choose the cost matrix
\begin{equation}\label{eq:chicost}
  c=
  \begin{pmatrix}
  0 & L\\
  L & 0
  \end{pmatrix},
  \qquad L>2\eps.
\end{equation}
Start from $f=(0,0)$.  For $y=1$, the update $g_1=T_\alpha f(1)$ solves
\begin{equation*}
  \frac12\left(1+\frac{g_1}{\eps}\right)_+
  +
  \frac12\left(1+\frac{g_1-L}{\eps}\right)_+
  =1.
\end{equation*}
Because $L>2\eps$, the solution is $g_1=\eps$, and only $x=1$ is active.
Similarly $g_2=\eps$, and only $x=2$ is active.  Hence the derivative kernels
are
\begin{equation*}
  P_1=\delta_1,
  \qquad
  P_2=\delta_2.
\end{equation*}
By Proposition~\ref{prop:derivative},
\begin{equation*}
  DT_\alpha(0)[h]=(-h_1,-h_2).
\end{equation*}
Applying the second block update gives the same derivative structure, hence
\begin{equation*}
  D(T_\beta T_\alpha)(0)=\id.
\end{equation*}
The local contraction factor of the double transform in variation seminorm is
therefore exactly one.

\begin{remark}
This example shows that strict contraction cannot be true for arbitrary
$\phi$-divergences without an overlap/interiority assumption.  It also explains
the mechanism: different target points see disjoint active source sets, so the
Markov kernels $P_y$ have zero common mass.
\end{remark}

\section{Noncompact and weighted extensions}\label{sec:noncompact}

The compact bounded-cost theory obtains strict contraction from a uniform
minorization: all derivative kernels share a fixed amount of mass.  On
noncompact spaces with unbounded costs, this is usually false globally.  There
are two natural ways to proceed.

\subsection{Weighted variation and Harris-type estimates}

Replace the oscillation seminorm by a weighted quotient seminorm, for example
\begin{equation}\label{eq:weightedosc}
  \|h\|_{V,\mathrm{var}}
  \defeq
  \inf_{a\in\R}
  \left\|\frac{h-a}{V}\right\|_\infty,
  \qquad V\geq 1.
\end{equation}
The derivative kernels $P_y^f$ then need not satisfy a global Doeblin
condition.  It is enough to prove a Harris-type combination of
\begin{enumerate}[label=(\roman*)]
  \item a minorization on a large set where the cost is controlled and
  $\Phi''$ is bounded away from zero;
  \item a drift estimate showing that $P_y^f$ does not place too much mass in
  the tails relative to $V$;
  \item a priori weighted bounds for the potentials and for
  $f\oplus g-c$.
\end{enumerate}
This is the natural general-$\phi$ analogue of the way KL Hilbert-metric
arguments have been adapted to unbounded kernels by modifying the cone and
controlling tail growth \cite{eckstein2025}.  The difference is that the kernel
to control is not the Gibbs kernel $e^{-c/\eps}$ itself, but the curvature
kernel $\Phi''(f\oplus g-c)$ normalized in each fiber.

\subsection{Weighted Carlier estimates}

The Carlier route asks for a different set of estimates.  One needs weighted
versions of
\begin{equation*}
  \mu_I\leq \Phi''\leq L_I
\end{equation*}
not necessarily on a bounded interval, but on the region reached by the
iterates.  A typical assumption would be that for some weights $V_i$,
\begin{equation*}
  |u_i^t(x_i)|\lesssim V_i(x_i),
  \qquad
  \sum_i u_i^t(x_i)-c(x)\in I(x)
\end{equation*}
with
\begin{equation*}
  0<\mu(x)\leq \Phi''(s)\leq L(x)<\infty,
  \qquad s\in I(x),
\end{equation*}
and with the weighted ratios integrable enough to reproduce the two scalar
estimates \eqref{eq:strongPhi}--\eqref{eq:smoothPhi} in a weighted $L^2$ norm.
This is plausible for divergences whose curvature behaves regularly in the
tails, but it is no longer automatic.  The key task is always the same:
prevent the iterates from approaching a region where $\Phi''$ either vanishes
or explodes too fast.

\section{Checklist for extending a KL proof to a given divergence}

Given a specific $\phi$, the following checklist isolates the work needed to
obtain strict contraction or linear convergence.

\begin{enumerate}[label=\textbf{Step \arabic*.},leftmargin=2.8cm]
  \item \textbf{Dual integrand.} Compute
  $\Phi(s)=\eps\phi_+^*(s/\eps)$, its derivative $\vartheta=\Phi'$, and its
  curvature $\Phi''$ on the interior of its domain.

  \item \textbf{Interior transform.} Show that the block equation
  \begin{equation*}
    \int \vartheta(f+\gamma-c)\dd\alpha=1
  \end{equation*}
  has a unique interior solution along the iterates.  If the solution can hit
  the boundary, strict contraction may fail.

  \item \textbf{A priori range.} Prove that $f\oplus g-c$ remains in a region
  $I$ or in a weighted region where the curvature is controlled.

  \item \textbf{Variation contraction route.} Verify a uniform overlap
  condition for the curvature kernels
  \begin{equation*}
    P_y^f(\dd x)\propto \Phi''(f(x)+T_\alpha f(y)-c(x,y))\alpha(\dd x).
  \end{equation*}
  A sufficient condition is $0<\inf_I\Phi''\leq \sup_I\Phi''<\infty$.

  \item \textbf{$L^2$ coordinate-descent route.} Verify the scalar inequalities
  \eqref{eq:strongPhi}--\eqref{eq:smoothPhi} on the same region and reproduce
  the sufficient-decrease and gap-to-step estimates.

  \item \textbf{Gauge and compactness.} Fix the additive constants and prove
  convergence of subsequences or compactness in the relevant topology.  Strict
  convexity on the quotient then identifies the limit.
\end{enumerate}

\section{Conclusion}

The order-theoretic part of Sinkhorn is robust.  For a broad class of
$\phi$-divergences, the generalized transform is order reversing and
translation covariant; the double transform is therefore non-expansive in the
variation seminorm.  Strict contraction is a different matter.  In KL it is
supported by the exponential factorization, which makes the algorithm a
projective positive-kernel iteration.  General $\phi$-regularization loses this
multiplicative structure.

The replacement principle is:
\begin{equation*}
  \boxed{\text{strict contraction comes from overlap of the curvature kernels}}
\end{equation*}
with kernels proportional to $\Phi''(f\oplus g-c)$.  If the iterates stay in an
interior strip where $0<\mu\leq \Phi''\leq L<\infty$, then the variation and
Carlier-style $L^2$ analyses extend with constants depending on $\mu/L$.  If
that interiority fails, sparse divergences can have disjoint active sets and
local contraction factor one.

\begin{thebibliography}{99}

\bibitem{beck2013}
A.~Beck and L.~Tetruashvili.
\newblock On the convergence of block coordinate descent type methods.
\newblock \emph{SIAM Journal on Optimization}, 23(4):2037--2060, 2013.

\bibitem{birkhoff1957}
G.~Birkhoff.
\newblock Extensions of Jentzsch's theorem.
\newblock \emph{Transactions of the American Mathematical Society},
85(1):219--227, 1957.

\bibitem{carlier2022}
G.~Carlier.
\newblock On the linear convergence of the multi-marginal Sinkhorn algorithm.
\newblock \emph{SIAM Journal on Optimization}, 32(2):786--794, 2022.
\newblock Preprint version: \url{https://mathtube.org/sites/default/files/lecture-extra-files/linear-sinkhorn.pdf}.

\bibitem{cuturi2013}
M.~Cuturi.
\newblock Sinkhorn distances: Lightspeed computation of optimal transport.
\newblock In \emph{Advances in Neural Information Processing Systems}, 2013.

\bibitem{dimarino2020}
S.~Di Marino and A.~Gerolin.
\newblock Optimal transport losses and Sinkhorn algorithm with general convex regularization.
\newblock \emph{arXiv:2007.00976}, 2020.

\bibitem{dobrushin1956}
R.~L. Dobrushin.
\newblock Central limit theorem for nonstationary Markov chains. I, II.
\newblock \emph{Theory of Probability and its Applications}, 1:65--80, 329--383, 1956.

\bibitem{eckstein2025}
S.~Eckstein.
\newblock Hilbert's projective metric for functions of bounded growth and exponential convergence of Sinkhorn's algorithm.
\newblock \emph{Probability Theory and Related Fields}, 2025.
\newblock See also \emph{arXiv:2311.04041}.

\bibitem{eveson1995}
S.~P. Eveson and R.~D. Nussbaum.
\newblock An elementary proof of the Birkhoff--Hopf theorem.
\newblock \emph{Mathematical Proceedings of the Cambridge Philosophical Society}, 117(1):31--55, 1995.

\bibitem{franklin1989}
J.~Franklin and J.~Lorenz.
\newblock On the scaling of multidimensional matrices.
\newblock \emph{Linear Algebra and its Applications}, 114/115:717--735, 1989.

\bibitem{ghosal2025}
P.~Ghosal and M.~Nutz.
\newblock On the convergence rate of Sinkhorn's algorithm.
\newblock Preprint, 2025.
\newblock \url{https://www.math.columbia.edu/~mnutz/docs/Sinkhorn_rate.pdf}.

\bibitem{nutz2022}
M.~Nutz.
\newblock Introduction to entropic optimal transport.
\newblock Lecture notes, 2022.
\newblock \url{https://www.math.columbia.edu/~mnutz/docs/EOT_lecture_notes.pdf}.

\bibitem{peyre2019}
G.~Peyre and M.~Cuturi.
\newblock Computational optimal transport.
\newblock \emph{Foundations and Trends in Machine Learning}, 11(5--6):355--607, 2019.

\bibitem{sinkhorn1967}
R.~Sinkhorn and P.~Knopp.
\newblock Concerning nonnegative matrices and doubly stochastic matrices.
\newblock \emph{Pacific Journal of Mathematics}, 21(2):343--348, 1967.

\bibitem{terjek2022}
D.~Terjek and D.~Gonzalez-Sanchez.
\newblock Optimal transport with $f$-divergence regularization and generalized Sinkhorn algorithm.
\newblock In \emph{Proceedings of AISTATS 2022}, PMLR 151:5135--5165, 2022.
\newblock \url{https://proceedings.mlr.press/v151/terjek22a.html}.

\bibitem{villani2009}
C.~Villani.
\newblock \emph{Optimal Transport: Old and New}.
\newblock Springer, 2009.

\end{thebibliography}

\end{document}
